\section{Figure 2: the orientability route}

This section separates the completed smooth-atlas proof for the torus from the
Klein bottle obstruction lemmas and the unfinished local homology route.
The future targets have no Lean declaration attached and are explicitly marked
not ready. Their dependency arrows describe the proposed proof route.

\subsection{Smooth oriented atlases}

\begin{definition}[Oriented atlas]\label{def:oriented-atlas}
  %\lean{GuibasStolfi.OrientedAtlas, GuibasStolfi.HasOrientedAtlas}

  An oriented atlas is a charted-space structure modeled on $\mathbb R^2$
  whose transitions are smooth and have positive Jacobian determinant at
  every point of every overlap. The charted-space axioms ensure that charts
  cover the whole space. Write $\operatorname{OA}(M)$ for existence of such an atlas.
\end{definition}

\begin{lemma}[Transport of oriented atlases]\label{lem:orientation-invariant}
  \uses{def:oriented-atlas, def:isomorphism-equivalent}
  %\lean{GuibasStolfi.OrientedAtlas.isManifold,
  %      GuibasStolfi.hasOrientedAtlas_iff_of_homeomorph,
  %      GuibasStolfi.Subdivision.Equivalent.hasOrientedAtlas_iff}

  An oriented atlas supplies a Mathlib smooth manifold structure. Existence
  of an oriented atlas is invariant under arbitrary homeomorphisms, and hence
  under subdivision equivalence.
\end{lemma}
\begin{proof}

  Smooth transitions give the manifold structure. For a homeomorphism $h$,
  transport each chart $e$ to $e\circ h^{-1}$. The transition maps are unchanged,
  since the two copies of $h$ cancel. Thus their smoothness and determinants
  are unchanged. No differentiability of $h$ is required.
\end{proof}

\begin{lemma}[Torus interval-lift atlas]\label{lem:torus-lift-atlas}
  \uses{def:torus-cells}
  %\lean{GuibasStolfi.Torus.orientationChartedSpace,
  %      GuibasStolfi.Torus.transition_locally_translation,
  %      GuibasStolfi.Torus.transition_det_eq_one}

  Lift the two circle coordinates into intervals $(a_i,a_i+1)$. These
  charts cover every torus point. Every overlap transition is locally a
  translation and has derivative the identity, with determinant $1$.
\end{lemma}
\begin{proof}

  For a representative $r\in[0,1)^2$, choose the cut $a=r-(1/2,1/2)$,
  which places $r$ inside its target rectangle. At a point $x$ of an overlap
  set $c=e'(q_T(x))-x$. The quotient of $c$ is zero. Openness of the target
  implies that $y+c$ lies there for all $y$ near $x$, and the inverse chart
  relation gives $e'(q_T(y))=y+c$. Differentiate this local translation.
\end{proof}

\begin{proposition}[Torus orientability]\label{prop:torus-orientable}
  \uses{lem:torus-lift-atlas, def:oriented-atlas}
  %\lean{GuibasStolfi.Torus.orientable, GuibasStolfi.Torus.standardTorus_orientable}

  The torus example admits an oriented atlas, as does Mathlib's unit
  two-dimensional torus.
\end{proposition}
\begin{proof}
  \uses{lem:orientation-invariant}

  The interval-lift atlas has all transition determinants equal to $1$.
  Transport it along the homeomorphism to Mathlib's torus.
\end{proof}

\subsection{Completed Klein bottle obstruction lemmas}

\begin{lemma}[The lifted gluing reverses plane orientation]\label{lem:klein-twist}
  \uses{def:klein-quotient}
  %\lean{GuibasStolfi.KleinBottle.quotientMap_twist,
  %      GuibasStolfi.KleinBottle.twist_det_eq_neg_one}

  For $g(x,y)=(x+1,-y)$ one has $q_K\circ g=q_K$, and $\det Dg=-1$.
\end{lemma}
\begin{proof}

  The quotient identifies a point with its $\tau$-image. Differentiating
  the plane lift gives $\operatorname{diag}(1,-1)$.
\end{proof}

\begin{lemma}[Sign-change obstruction]\label{lem:klein-sign-obstruction}
  \uses{def:klein-quotient}
  %\lean{GuibasStolfi.KleinBottle.exists_zero_of_flip_antiinvariant,
  %      GuibasStolfi.KleinBottle.no_nonzero_flip_antiinvariant}

  Every continuous $f:C\to\mathbb R$ satisfying $f(\tau(p))=-f(p)$
  vanishes somewhere. In particular, such a function cannot be everywhere nonzero.
\end{lemma}
\begin{proof}

  The torus $C$ is connected. Values at $p$ and $\tau(p)$ have opposite signs,
  so the intermediate value theorem gives a zero.
\end{proof}

\begin{lemma}[No continuous quotient section]\label{lem:klein-no-section}
  \uses{def:klein-quotient}
  %\lean{GuibasStolfi.KleinBottle.projection_no_continuous_section}

  There is no continuous $s:K\to C$ satisfying $\pi\circ s=\mathrm{id}_K$.
\end{lemma}
\begin{proof}

  If $s$ existed, the closed sets $\{p:s(\pi(p))=p\}$ and
  $\{p:s(\pi(p))=\tau(p)\}$ would cover $C$. They are disjoint since
  $\tau$ has no fixed points, and both are nonempty by evaluating at $s(z)$
  and $\tau(s(z))$. This contradicts connectedness of $C$.
\end{proof}
These lemmas do not yet imply that every possible oriented atlas on $K$ is
impossible. The missing implication is that any such atlas would supply a
section of this particular quotient. A negative determinant in chosen
coordinates alone does not establish that implication.

\subsection{Coordinate-independent local homology}

\begin{definition}[Relative singular integral homology]\label{def:relative-homology}
  %\lean{GuibasStolfi.LocalOrientation.singularPairFunctor,
  %      GuibasStolfi.LocalOrientation.relativeHomologyFunctor}

  Apply the singular simplicial set functor to a topological pair $(M,A)$,
  then take Mathlib's relative simplicial homology with coefficient object
  $\mathbb Z$ in the category of abelian groups. Denote it $H_n(M,A;\mathbb Z)$.
  The relative chain complex is the cokernel of the subspace's chain inclusion.
\end{definition}

\begin{definition}[Local orientations]\label{def:local-orientation}
  \uses{def:relative-homology}
  %\lean{GuibasStolfi.LocalOrientation.localHomology, GuibasStolfi.LocalOrientation.At}

  Put $L_x(M)=H_2(M,M\setminus\{x\};\mathbb Z)$.
  A local orientation at $x$ is an additive equivalence
  $o_x:\mathbb Z\to L_x(M)$, specifying the generator $o_x(1)$.
\end{definition}

\begin{lemma}[Exactly two generators, conditionally]\label{lem:two-local-generators}
  \uses{def:local-orientation}
  %\lean{GuibasStolfi.OrientationGenerator.eq_or_eq_reverse,
  %      GuibasStolfi.OrientationGenerator.signEquiv,
  %      GuibasStolfi.OrientationGenerator.reverse_ne}

  Given one additive equivalence $o:\mathbb Z\to G$, all such equivalences
  are exactly $o$ and $n\mapsto o(-n)$, and these are distinct.
\end{lemma}
\begin{proof}

  The comparison of two equivalences is an additive automorphism of
  $\mathbb Z$, determined by its value at $1$. Its inverse shows this value
  is a unit, hence $1$ or $-1$. Injectivity distinguishes the two signs.
\end{proof}
This does not prove that any given surface point has a local generator.
The plane computation required for that conclusion is still open in this development.

\begin{definition}[Coherent homological orientation]\label{def:coherent-orientation}
  \uses{def:local-orientation}
  %\lean{GuibasStolfi.LocalOrientation.toLocal,
  %      GuibasStolfi.LocalOrientation.IsCoherent, GuibasStolfi.LocalOrientation.Coherent}

  For $x\in U$, complement inclusion gives a map
  $\rho_{U,x}:H_2(M,M\setminus U;\mathbb Z)\to L_x(M)$.
  A field $(o_x)$ of local orientations is coherent if every point has an
  open neighborhood $U$ and one class $c\in H_2(M,M\setminus U;\mathbb Z)$
  with $\rho_{U,y}(c)=o_y(1)$ for all $y\in U$.
\end{definition}

\begin{lemma}[Homeomorphism invariance of coherence]\label{lem:homological-orientation-invariant}
  \uses{def:coherent-orientation, def:isomorphism-equivalent}
  %\lean{GuibasStolfi.LocalOrientation.relativeHomeomorph_natural,
  %      GuibasStolfi.LocalOrientation.nonempty_coherent_iff,
  %      GuibasStolfi.LocalOrientation.isCoherent_reverse,
  %      GuibasStolfi.Subdivision.Equivalent.nonempty_coherent_iff}

  Arbitrary homeomorphisms transport coherent homological orientations.
  Existence is therefore invariant under subdivision equivalence. Reversing
  every local orientation preserves coherence.
\end{lemma}
\begin{proof}

  Transport the punctured pairs and representing neighborhood classes through
  the homeomorphism. Naturality with complement inclusions preserves the
  representing equations. For reversal negate the representing class.
\end{proof}

\subsection{Remaining targets for the orientation-cover proof}

The following standard mathematical steps need substantial formal development.
No excision axiom, local homology computation, or atlas-comparison axiom has
been introduced. The two orientation criteria above are not yet linked in Lean.
The current homology declarations use small Lean types (universe zero).

\begin{lemma}[Local excision]\label{lem:local-excision}
  \notready
  \uses{def:relative-homology}
  For a Hausdorff space $M$, an open neighborhood $U$ of $x$, and all $n$,
  inclusion induces an isomorphism
  $H_n(U,U\setminus\{x\};\mathbb Z)\cong H_n(M,M\setminus\{x\};\mathbb Z)$.
\end{lemma}
Excision for the relative singular homology construction was not found in
the pinned Mathlib checkout. The stronger general excision theorem would
give this neighborhood statement.

\begin{lemma}[Plane local homology]\label{lem:plane-local-homology}
  \notready
  \uses{def:local-orientation}
  The local group $L_x(\mathbb R^2)$ is infinite cyclic, with generators
  compatible under restriction to sufficiently small disk neighborhoods.
\end{lemma}

\begin{proposition}[Orientation double cover]\label{prop:orientation-cover}
  \notready
  \uses{lem:local-excision, lem:plane-local-homology, lem:two-local-generators,
        def:coherent-orientation}
  For a Hausdorff surface, the local orientations assemble into a topological
  double cover of the surface, and coherent fields give continuous sections.
\end{proposition}

\begin{lemma}[Local homology and Jacobian sign]\label{lem:jacobian-homology}
  \notready
  \uses{lem:plane-local-homology, lem:local-excision}
  Relative to the standard plane orientation, a smooth local diffeomorphism
  acts on local homology with the sign of its Jacobian determinant.
  In particular a plane reflection acts by negation.
\end{lemma}

\begin{proposition}[Oriented atlases give sections]\label{prop:atlas-gives-section}
  \notready
  \uses{def:oriented-atlas, prop:orientation-cover, lem:jacobian-homology}
  Every oriented atlas induces a coherent homological orientation and a
  continuous section of the orientation double cover.
\end{proposition}

\begin{proposition}[The Klein bottle orientation cover]\label{prop:klein-orientation-cover}
  \notready
  \uses{prop:orientation-cover, lem:jacobian-homology, lem:klein-twist,
        lem:klein-charts}
  The orientation double cover of $K$ is isomorphic, over $K$, to
  $\pi:C\to K$ from the Klein bottle construction.
\end{proposition}

\begin{proposition}[Klein bottle nonorientability]\label{prop:klein-nonorientable}
  \notready
  \uses{prop:atlas-gives-section, prop:klein-orientation-cover, lem:klein-no-section}
  The Klein bottle $K$ admits no oriented atlas.
\end{proposition}
The planned contradiction uses the proved no-section lemma after the two
preceding comparisons have been established.

\begin{proposition}[Nonequivalent subdivisions]\label{prop:torus-klein-nonequivalent}
  \notready
  \uses{prop:torus-orientable, prop:klein-nonorientable, lem:orientation-invariant,
        prop:torus-subdivision, prop:klein-subdivision}
  The constructed torus and Klein bottle subdivisions are not equivalent.
\end{proposition}
Together with their proved graph isomorphism, this would complete
Proposition~\ref{prop:graph-not-determining}. That final conclusion is still open.
